\documentclass[a4paper,12pt]{article}
\usepackage[T2A]{fontenc}
\usepackage[utf8]{inputenc}
\usepackage[russian]{babel}
\usepackage{amsmath,amssymb,mathtools}
\usepackage{geometry}
\geometry{left=22mm,right=22mm,top=25mm,bottom=23mm}
\usepackage{xcolor}
\usepackage{tikz}
\usetikzlibrary{arrows.meta,positioning,shapes.geometric,calc}
\usepackage{booktabs,tabularx}
\usepackage{enumitem}
\usepackage{fancyhdr}
\usepackage{titlesec}
\usepackage{microtype}
\usepackage{icomma}
\usepackage{hyperref}
\usepackage{setspace}
\definecolor{accent}{HTML}{008080}
\definecolor{darkaccent}{HTML}{005B5B}
\definecolor{ink}{HTML}{303639}
\definecolor{muted}{HTML}{656B6B}
\hypersetup{colorlinks=true,linkcolor=darkaccent,urlcolor=darkaccent,pdftitle={Чек-лист. Задачи прикладного характера},pdfauthor={Лёвин Артём Александрович}}
\setlength{\parindent}{0pt}
\setlength{\parskip}{5.5pt}
\linespread{1.16}
\color{ink}
\titleformat{\section}{\color{darkaccent}\Large\normalfont}{\thesection.}{0.7em}{}[\vspace{2pt}\color{accent}\titlerule]
\titlespacing*{\section}{0pt}{18pt}{13pt}
\titleformat{\subsection}{\color{darkaccent}\normalsize\bfseries}{\thesubsection}{0.7em}{}
\titlespacing*{\subsection}{0pt}{11pt}{5pt}
\setlist[enumerate]{leftmargin=1.7em,itemsep=8pt,topsep=7pt}
\pagestyle{fancy}
\fancyhf{}
\fancyhead[L]{\small Задачи прикладного характера}
\fancyhead[R]{\small Анна Баннова}
\fancyfoot[C]{\small\thepage}
\renewcommand{\headrulewidth}{0.4pt}
\renewcommand{\headrule}{\hbox to\headwidth{\color{accent}\leaders\hrule height \headrulewidth\hfill}}
\setlength{\headheight}{25pt}
\newcommand{\idea}[1]{\par\textcolor{darkaccent}{\textbf{Ключевая идея.}} #1\par}
\newcommand{\warn}[1]{\par\textcolor{darkaccent}{\textbf{Проверьте себя.}} #1\par}
\newcommand{\ex}[1]{\par\vspace{4pt}{\color{darkaccent}\textbf{#1}}\par}
\begin{document}
\thispagestyle{empty}
\begin{center}
\vspace*{0.09\textheight}
{\large\color{darkaccent} МАТЕМАТИКА}\par
\vspace{1.10cm}
{\Huge Чек-лист}\par
\vspace{0.90cm}
{\LARGE\color{accent} Задачи прикладного характера}\par
\vspace{0.36cm}
{\large Тепловое расширение, прибыль, изменение уровня воды}\par
\vspace{1.00cm}
{\large Подготовка к ЕГЭ по профильной математике}\par
\vfill
{\large Анна Баннова}\par
\vspace{0.55cm}
{\large 26 октября 2026 г.}\par
\vspace{0.95cm}
{\normalsize Лёвин Артём Александрович}\par
{\small эксклюзивно для Анны Банновой}\par
\vspace{1.05cm}
\begin{tikzpicture}[x=1cm,y=1cm]
\draw[darkaccent,line width=1.0pt] (0,0) .. controls (1.65,0.80) and (3.1,-0.75) .. (4.5,0)
.. controls (6.0,0.68) and (7.2,-0.62) .. (8.4,0.10);
\fill[darkaccent] (4.5,0) circle (1.5pt);
\end{tikzpicture}
\end{center}
\clearpage

\section{Как переводить текст задачи в математику}
В прикладной задаче сначала определите, \emph{какую величину нужно найти}, какие данные известны и как они связаны формулой. Особое внимание уделяйте словам об изменении величины.

\idea{Если величина выросла от начального значения $x_1$ до конечного $x_2$, её прирост равен $\Delta x=x_2-x_1$. Если требуется величина уменьшения, вычтите конечное значение из начального.}

\begin{center}
\begin{tikzpicture}[>=Latex,x=1cm,y=1cm]
\draw[line width=1pt,muted] (0,0) -- (6.5,0);
\draw[line width=2.5pt,accent] (0,0.35) -- (4.0,0.35);
\draw[line width=2.5pt,darkaccent] (0,0.80) -- (6.0,0.80);
\draw[<->,accent] (4.02,1.18) -- (6.0,1.18);
\node[above,accent] at (5.02,1.18) {$\Delta x$};
\node[right] at (6.0,0.80) {$x_2$};
\node[right] at (4.0,0.35) {$x_1$};
\end{tikzpicture}
\end{center}

\subsection{Размерности и степени десяти}
Перед подстановкой величины должны быть выражены в согласованных единицах:
\[
1\ \text{мм}=10^{-3}\ \text{м},\qquad 4{,}5\ \text{мм}=4{,}5\cdot10^{-3}\ \text{м}
=45\cdot10^{-4}\ \text{м}.
\]
Вычисляя, используйте свойства степеней:
\[
10^m\cdot10^n=10^{m+n},\qquad \frac{10^m}{10^n}=10^{m-n}.
\]
Например, $1{,}2\cdot10^{-5}=12\cdot10^{-6}$; при переносе запятой на один знак вправо показатель десятки уменьшается на единицу.

\subsection{Преобразование формул}
Для выражения $a(b+c)$ выполняется правило
\[
a(b+c)=ab+ac.
\]
Множитель $a$ умножает \emph{каждое слагаемое ровно один раз}. Произведение $bc$ внутри слагаемого рассматривается как единое целое.

\warn{Сначала преобразуйте общую формулу, затем подставляйте числа. Это помогает увидеть одинаковые выражения и сократить число действий.}

\section{Тепловое расширение рельса}
При нагревании от $0^\circ\text{C}$ до температуры $T$ длина рельса меняется по закону
\[
l(T)=l_0(1+\alpha T),
\]
где $l_0$ --- начальная длина, $l(T)$ --- длина при температуре $T$, $\alpha$ --- коэффициент теплового расширения.
Раскроем скобки и выразим удлинение:
\[
\begin{aligned}
l(T)&=l_0+l_0\alpha T,\\
\underbrace{l(T)-l_0}_{\Delta l}&=l_0\alpha T,\\
T&=\frac{\Delta l}{l_0\alpha}\quad(l_0\alpha\ne0).
\end{aligned}
\]

\ex{Пример 1. На сколько градусов нужно нагреть рельс?}
Дано: $l_0=10$ м, $\alpha=1{,}2\cdot10^{-5}\,{}^\circ\text{C}^{-1}$, удлинение $\Delta l=4{,}5$ мм. Требуется найти $T$.
\[
\begin{aligned}
\Delta l&=4{,}5\cdot10^{-3}=45\cdot10^{-4}\ \text{м},\\
\alpha&=12\cdot10^{-6}\,{}^\circ\text{C}^{-1},\\
T&=\frac{45\cdot10^{-4}}{10\cdot12\cdot10^{-6}}
=\frac{45}{120}\cdot10^2
=37{,}5^\circ\text{C}.
\end{aligned}
\]
\textbf{Проверка:} $10\cdot1{,}2\cdot10^{-5}\cdot37{,}5=0{,}0045\ \text{м}=4{,}5\ \text{мм}$.

\section{Минимальный объём производства}
Пусть компания продаёт одну единицу товара за $p$ рублей, затраты на одну единицу равны $v$ рублей, постоянные расходы за месяц --- $f$ рублей. Операционная прибыль при выпуске $q$ единиц:
\[
\pi(q)=q(p-v)-f.
\]
Если прибыль должна быть \emph{не менее} $A$ рублей, пишем $\pi(q)\ge A$.
При $p>v$:
\[
q(p-v)-f\ge A
\quad\Longleftrightarrow\quad
q\ge\frac{A+f}{p-v}.
\]
Поскольку число единиц товара целое, при необходимости округляем найденную нижнюю границу \emph{вверх}.

\ex{Пример 2. Каков наименьший месячный объём выпуска?}
Дано: $p=600$ руб., $v=300$ руб., $f=700\,000$ руб., требуемая прибыль $A=500\,000$ руб.
\[
\begin{aligned}
q(600-300)-700\,000&\ge500\,000,\\
300q&\ge1\,200\,000,\\
q&\ge4\,000.
\end{aligned}
\]
\textbf{Ответ:} наименьший объём --- $4\,000$ единиц. При таком выпуске прибыль равна ровно $500\,000$ руб. При $q=3\,999$ она составит только $499\,700$ руб.

\warn{Для поиска порога можно решить уравнение $\pi(q)=A$, но условие «не менее» означает неравенство. Если граница дробная, нужен ближайший \emph{больший} целый выпуск.}

\section{Подъём воды в колодце}
Для свободного падения в условии задачи используется зависимость пройденного расстояния от времени
\[
h=5t^2,
\]
где $h$ измеряется в метрах, $t$ --- в секундах. До дождя камень летел до воды дольше, после подъёма уровня воды --- меньше.

\begin{center}
\begin{tikzpicture}[x=1cm,y=1cm]
% walls
\draw[muted,line width=0.8pt] (0,0) -- (0,2.4) -- (1.4,2.4) -- (1.4,0);
\draw[muted,line width=0.8pt] (3.2,0) -- (3.2,2.4) -- (4.6,2.4) -- (4.6,0);
\draw[accent,line width=1.1pt] (0,0.4)--(1.4,0.4);
\draw[accent,line width=1.1pt] (3.2,1.15)--(4.6,1.15);
\draw[<->,>=Latex,darkaccent] (0.65,0.45)--(0.65,2.35);
\draw[<->,>=Latex,darkaccent] (3.85,1.2)--(3.85,2.35);
\draw[<->,>=Latex,accent] (5.0,0.4)--(5.0,1.15);
\node[right] at (5.07,0.77) {$\Delta h$};
\node[below] at (0.7,-0.02) {до дождя};
\node[below] at (3.9,-0.02) {после дождя};
\node[left] at (0.58,1.42) {$h_1$};
\node[left] at (3.78,1.75) {$h_2$};
\end{tikzpicture}
\end{center}

Пусть время падения до дождя равно $t_1$, после дождя --- $t_2$. При подъёме воды $t_2<t_1$, а высота подъёма воды равна разности прежнего и нового расстояний:
\[
\Delta h=h_1-h_2=5t_1^2-5t_2^2=5(t_1-t_2)(t_1+t_2).
\]
Последний переход основан на формуле разности квадратов $a^2-b^2=(a-b)(a+b)$.

\ex{Пример 3. Насколько поднялась вода?}
До дождя камень падал $0{,}6$ с. После дождя время уменьшилось на $0{,}2$ с, значит, $t_2=0{,}6-0{,}2=0{,}4$ с.
\[
\begin{aligned}
\Delta h&=5(0{,}6^2-0{,}4^2)
=5(0{,}6-0{,}4)(0{,}6+0{,}4)\\
&=5\cdot0{,}2\cdot1=1\ \text{м}.
\end{aligned}
\]
Проверка по отдельным расстояниям: $h_1=1{,}8$ м, $h_2=0{,}8$ м; их разность равна $1$ м.

\subsection*{Проверка вычислений}
При работе со степенями отдельно считайте числовые коэффициенты и отдельно степени десяти. Например, $6/12=0{,}5$, затем учитывается множитель $10^k$; так проще избежать ошибки в положении запятой.

Для задачи о прибыли проверьте найденный целый выпуск и выпуск на одну единицу меньше. Для задачи о подъёме воды вычислите оба расстояния по формуле $h=5t^2$ и сравните их разность с ответом.

\clearpage
\thispagestyle{fancy}
\begin{center}
{\Large\color{darkaccent} Алгоритм: от текста к ответу}\par
\vspace{0.60cm}
\end{center}
\begin{center}
\tikzset{
 stage/.style={draw=accent!70,rounded corners=2pt,minimum width=10.3cm,text width=9.6cm,align=center,minimum height=1.15cm,inner sep=7pt,font=\normalsize},
 choice/.style={draw=accent!80,diamond,aspect=2.7,text width=4.4cm,align=center,inner sep=2pt,font=\small},
 arr/.style={-Latex,very thick,draw=darkaccent}
}
\begin{tikzpicture}[node distance=0.61cm and 1.2cm]
\node[stage] (a) {Прочитайте условие. Что известно и что требуется найти?};
\node[stage,below=of a] (b) {Обозначьте величины; при изменении сделайте схему «до / после».};
\node[stage,below=of b] (c) {Приведите единицы к одной системе. Переведите слова в формулу или неравенство.};
\node[stage,below=of c] (d) {Раскройте скобки, вынесите общий множитель, примените свойства степеней.};
\node[stage,below=of d] (e) {Выразите искомую величину и вычислите.};
\node[choice,below=0.80cm of e] (f) {Нужен наименьший целый выпуск?};
\node[stage,below left=1.05cm and -0.75cm of f,minimum width=5.1cm,text width=4.6cm] (g) {Да: округлите нижнюю границу вверх.};
\node[stage,below right=1.05cm and -0.75cm of f,minimum width=5.1cm,text width=4.6cm] (h) {Нет: сохраните вычисленное значение.};
\node[stage,below=1.06cm of f,yshift=-2.55cm] (i) {Проверьте подстановкой, размерностью и смыслом ответа.};
\draw[arr] (a)--(b);
\draw[arr] (b)--(c);
\draw[arr] (c)--(d);
\draw[arr] (d)--(e);
\draw[arr] (e)--(f);
\draw[arr] (f.south west) -| node[near start,above,sloped,font=\small]{да} (g.north);
\draw[arr] (f.south east) -| node[near start,above,sloped,font=\small]{нет} (h.north);
\draw[arr] (g.south) |- (i.west);
\draw[arr] (h.south) |- (i.east);
\end{tikzpicture}
\end{center}
\clearpage

\section*{Тренировочный блок. Путь А}
\addcontentsline{toc}{section}{Тренировочный блок}
\textcolor{muted}{10 задач --- от стандартных подстановок к задачам с дополнительной проверкой. Решайте самостоятельно; ответы приведены на следующей странице.}
\vspace{0.10cm}
\begin{enumerate}[label=\textbf{\arabic*.},leftmargin=1.55em,itemsep=13pt,parsep=0pt]
\item Рельс при $0^\circ\text{C}$ имеет длину $8$ м, коэффициент расширения $2\cdot10^{-5}\,{}^\circ\text{C}^{-1}$. При какой температуре рельс удлинится на $3{,}2$ мм? Используйте $l(T)=l_0(1+\alpha T)$.
\item При $0^\circ\text{C}$ длина рельса $12$ м, $\alpha=1{,}5\cdot10^{-5}\,{}^\circ\text{C}^{-1}$. Найдите температуру, при которой его длина увеличится на $5{,}4$ мм.
\item Длина рельса при $0^\circ\text{C}$ равна $15$ м, $\alpha=1{,}2\cdot10^{-5}\,{}^\circ\text{C}^{-1}$. На сколько миллиметров удлинится рельс при нагревании до $40^\circ\text{C}$?
\item Товар продают за $800$ руб.; затраты на единицу --- $500$ руб., постоянные расходы --- $450\,000$ руб. Каков наименьший месячный выпуск для прибыли не менее $750\,000$ руб.?
\item Цена товара --- $950$ руб., переменные затраты на единицу --- $650$ руб., постоянные расходы --- $360\,000$ руб. Найдите минимальный месячный выпуск для прибыли не менее $840\,000$ руб.
\item Цена --- $410$ руб., переменные затраты --- $260$ руб. на единицу, постоянные расходы --- $210\,000$ руб. Определите наименьший \emph{целый} выпуск для прибыли не менее $500\,000$ руб.
\item Камень до дождя падал до воды $0{,}8$ с, после дождя время уменьшилось на $0{,}2$ с. На сколько метров поднялся уровень воды? Используйте $h=5t^2$.
\item До дождя камень падал до воды $1{,}0$ с, после дождя --- на $0{,}4$ с меньше. На сколько метров поднялась вода? $h=5t^2$.
\item Время падения камня до дождя составляло $0{,}9$ с и уменьшилось после дождя на $0{,}3$ с. Найдите подъём воды в метрах. $h=5t^2$.
\item Время падения до дождя было $1{,}2$ с. После подъёма воды оно стало меньше на $0{,}25$ с. Найдите, на сколько метров поднялась вода. $h=5t^2$.
\end{enumerate}
\clearpage
\section*{Ответы к тренировочному блоку}
\vspace{0.25cm}
\renewcommand{\arraystretch}{1.42}
\begin{center}
\begin{tabular}{@{}c p{5.9cm} l@{}}
\toprule
№ & Тема & Ответ \\
\midrule
1 & Тепловое расширение & $20^\circ\text{C}$\\
2 & Тепловое расширение & $30^\circ\text{C}$\\
3 & Тепловое расширение & $7{,}2$ мм\\
4 & Производство и прибыль & $4\,000$ ед.\\
5 & Производство и прибыль & $4\,000$ ед.\\
6 & Производство и прибыль & $4\,734$ ед.\\
7 & Подъём уровня воды & $1{,}4$ м\\
8 & Подъём уровня воды & $3{,}2$ м\\
9 & Подъём уровня воды & $2{,}25$ м\\
10 & Подъём уровня воды & $2{,}6875$ м\\
\bottomrule
\end{tabular}
\end{center}
\vspace{0.90cm}
{\color{darkaccent}\large Итоговая памятка}\par
\vspace{0.20cm}
\renewcommand{\arraystretch}{1.45}
\begin{tabularx}{\linewidth}{@{}p{0.34\linewidth}X@{}}
\toprule
Сигнал в условии & Действие \\
\midrule
«Увеличился на\ldots» & Найдите разность конечной и начальной величин.\\
«При температуре\ldots» & Подставьте изменение в формулу $\Delta l=l_0\alpha T$.\\
«Прибыль не менее\ldots» & Составьте неравенство; найдите минимальное целое $q$.\\
«Время стало меньше» & Сначала найдите новое время $t_2$, затем разность $5t_1^2-5t_2^2$.\\
В формуле есть скобки или квадраты & Проверьте распределительный закон, вынесение множителя, разность квадратов.\\
\bottomrule
\end{tabularx}
\vfill
\begin{center}
{\small\color{muted}Лёвин Артём Александрович · эксклюзивно для Анны Банновой}
\end{center}
\end{document}
